% Separate analytical increment; do not merge without explicit user approval.
% Label prefix: inc:SGWR:v1:
\documentclass[11pt]{article}
\usepackage[T1]{fontenc}
\usepackage[utf8]{inputenc}
\usepackage{lmodern,amsmath,amssymb,amsthm,mathtools,microtype}
\usepackage[margin=0.95in]{geometry}
\usepackage{xurl,booktabs,array}
\usepackage[hidelinks]{hyperref}
\usepackage{fancyhdr}
\pagestyle{fancy}\fancyhf{}
\fancyhead[L]{\small Signed growth and weak-side retention}
\fancyhead[R]{\small Separate review increment v1}
\fancyfoot[C]{\thepage}
\setlength{\headheight}{14pt}
\setlength{\emergencystretch}{3em}
\allowdisplaybreaks[2]
\numberwithin{equation}{section}
\newtheorem{theorem}{Theorem}[section]
\newtheorem{proposition}[theorem]{Proposition}
\newtheorem{lemma}[theorem]{Lemma}
\newtheorem{corollary}[theorem]{Corollary}
\theoremstyle{remark}\newtheorem{remark}[theorem]{Remark}
\newcommand{\E}{\mathbb E}
\newcommand{\R}{\mathbb R}
\newcommand{\dd}{\,d}
\newcommand{\norm}[1]{\left\lVert#1\right\rVert}
\newcommand{\1}{\mathbf1}
\title{Signed growth, weak-coordinate amplification,\\and a conditional retention--passage interface}
\author{Separate analytical increment for independent review}
\date{Version 1 --- October 3, 2026}
\begin{document}\maketitle
\begin{abstract}
We replace loss-only retention bookkeeping by a signed-growth estimate.
For the exact balanced population, a weak-coordinate amplitude converts the
geometric factor in radial growth into an endpoint factor. This identity does
not require input/output alignment after initialization. An exact Gaussian
sector estimate gives a sufficient bound on its residual-matrix defect.
We also prove a uniform contraction estimate for the frozen resident weak
system on the AN06 ratio interval. Finally, we verify the exponent algebra
for the global profile multiplier quoted by the user. The cited ancestry/tail
file was not accessible, so that multiplier is an explicit conditional
premise, not a theorem verified here. None of these results proves the
initialized angular guards, weak capture, full two-sided passage, or
finite-noise Theorems A and B.
\end{abstract}

\paragraph{Provenance and review contract.}
The exact original-flow equations and balance are imported from
\path{RELU_POPULATION_FOUNDATIONS_v2.tex}, P1--P2; global finite-time
regularity and positive mass are imported from P5. The earlier negative-budget
retention bound is in \path{LIMIT_LEARNED_PERSISTENCE_RETENTION_v1.tex},
Proposition 5.1. The resident root is \texttt{pa:weakroot} in
\path{THEOREM_A_ANALYTICAL_PROGRESS.tex}; the compact ratio bounds are from
\path{AN06_learned_orbit_witnesses_v2.tex}.
The user quotes Theorem 5.1 of
\path{AN02_strong_ancestry_and_tail_transport_v1.tex}, labels
\texttt{inc:AN02:tail:v1:profiletheorem} and
\texttt{inc:AN02:tail:v1:globalG}. That file is unavailable in this review.
Its hypotheses, definition of $b_c$, reference characteristic, and exact
one-sided conclusion are not reconstructed here.
All new assertions are displayed with their assumptions and proofs. No saved
trajectory, floating-point sign test, or interval certificate is a premise.
No earlier file has been altered, and this increment is not approved for an
automatic merge.

\section{Signed growth for a fixed label set}
Use normalized time $\tau=\mu_1^2t_{\rm phys}$ and
\[
 P_1=\mathcal N(e_1,q^2I_2),\quad
 P_2=\mathcal N(\rho e_2,q^2I_2),\quad P=(P_1+P_2)/2,
 \qquad \lambda=\rho^2.
\]
The label law is $\lambda_0(d\alpha)=d\alpha/(2\pi)$.
Write $U_\alpha=\sqrt{m_\alpha}a_\alpha$,
$W_\alpha=\sqrt{m_\alpha}b_\alpha$, with unit $a_\alpha,b_\alpha$.
All output contributions are retained in
\[
 R_\tau(a)=\E_P[(X-f_\tau(X))X^\top\1_{\{a^\top X>0\}}].
\]
The exact equations are
\begin{equation}\label{inc:SGWR:v1:flow}
 U_\alpha'=R_\tau(a_\alpha)^\top W_\alpha,
 \quad W_\alpha'=R_\tau(a_\alpha)U_\alpha,
 \quad (\log m_\alpha)'=r_\alpha:=2b_\alpha^\top R_\tau(a_\alpha)a_\alpha.
\end{equation}
Let $C$ be a fixed measurable initial-label set, and set
$N_C(\tau)=\int_Cm_\alpha(\tau)\dd\lambda_0$.

\begin{proposition}[Signed Jensen lower bound]\label{inc:SGWR:v1:jensen}
Fix finite $\tau_a<\tau_b$ with $N_C(\tau_a)>0$, and suppose the integrals
below are absolutely integrable with the indicated weights. Define
\[
 I_\alpha=\int_{\tau_a}^{\tau_b}r_\alpha(t)\dd t,
 \qquad d\pi_a(\alpha)=\frac{m_\alpha(\tau_a)\1_C\dd\lambda_0}{N_C(\tau_a)}.
\]
Then
\begin{equation}\label{inc:SGWR:v1:signedjensen}
 \boxed{N_C(\tau_b)\ge N_C(\tau_a)
                  \exp\!\left(\int_C I_\alpha\dd\pi_a\right).}
\end{equation}
If $\int I_\alpha\dd\pi_a\ge
\lambda(\tau_b-\tau_a)-B_0-b\log(1/S_0)$ with $0<S_0<1$, then
\[
 N_C(\tau_b)\ge e^{-B_0}N_C(\tau_a)
                 e^{\lambda(\tau_b-\tau_a)}S_0^b.
\]
There is also an exact current-weight formula:
\begin{equation}\label{inc:SGWR:v1:currentweights}
 \log\frac{N_C(\tau_b)}{N_C(\tau_a)}
 =\int_{\tau_a}^{\tau_b}\int_C r_\alpha(t)\dd\pi_t(\alpha)\dd t,
 \qquad d\pi_t=\frac{m_\alpha(t)\1_C\dd\lambda_0}{N_C(t)}.
\end{equation}
\end{proposition}
\begin{proof}
Integrating \eqref{inc:SGWR:v1:flow} gives
$m_\alpha(\tau_b)=m_\alpha(\tau_a)e^{I_\alpha}$. Integrate over $C$ and
apply convexity of the exponential to the probability measure $\pi_a$.
The second assertion is substitution. For \eqref{inc:SGWR:v1:currentweights},
differentiate the fixed-label integral: $N_C'=\int_Cm_\alpha r_\alpha
\dd\lambda_0$. Positive mass permits division by $N_C$. The stated
integrability assumptions, or the finite-horizon bounds of the foundations,
justify differentiation, Fubini, and integration of the logarithmic derivative.
\end{proof}

\begin{remark}
The original negative-budget inequality remains a valid weaker corollary,
obtained by dropping all positive growth. It is useful after a seed has been
constructed, but cannot recover amplification that it has discarded. The
initial-weight signed average and the time-dependent current-weight average
are different objects: the former supplies an inequality, the latter an
identity. For a later captured subset $C_{\rm cap}\subset C$, apply the
proposition to that fixed subset and retain its initial mass fraction.
Growth of the full reservoir does not prove growth of the captured part.
\end{remark}

\section{An exact weak-coordinate amplitude}
The radial approximation $r\simeq\lambda\sin^2\theta$ need not be bounded
by a fixed angular minimum throughout the interval. A coordinate amplitude
retains the associated alignment gain exactly.
Define
\begin{equation}\label{inc:SGWR:v1:amplitude}
 z_\alpha=\frac{U_{\alpha,2}+W_{\alpha,2}}2,
 \qquad h_\alpha=z_\alpha^2,
 \qquad \chi_\alpha=\frac{a_{\alpha,2}+b_{\alpha,2}}2.
\end{equation}
Balance gives $h_\alpha=m_\alpha\chi_\alpha^2\le m_\alpha$.
Choose an absolutely continuous or measurable scalar $c(\tau)\in[0,1]$,
and decompose, exactly,
\begin{equation}\label{inc:SGWR:v1:defect}
 R_\tau(a_\alpha)
 =\frac\lambda2(1-c(\tau))e_2e_2^\top+\mathcal B_\alpha(\tau).
\end{equation}
The scalar $c$ is a comparison response, not automatically the mass of a
family. Choosing $c=N$ requires the residual estimate that makes that choice
useful; the identities do not assume that it is exact.

\begin{proposition}[Signed weak-amplitude equation]\label{inc:SGWR:v1:ampflow}
On an interval on which $\chi_\alpha>0$,
\begin{align}
 (\log h_\alpha)'&=\lambda(1-c)+\varepsilon_\alpha,
 \label{inc:SGWR:v1:amprate}\\
 \varepsilon_\alpha
 &=\frac{(\mathcal B_\alpha^\top W_\alpha)_2
                    +(\mathcal B_\alpha U_\alpha)_2}{z_\alpha}
 =2(\mathcal B_\alpha)_{22}
 +\frac{(\mathcal B_\alpha)_{12}W_{\alpha,1}
           +(\mathcal B_\alpha)_{21}U_{\alpha,1}}{z_\alpha}.
 \nonumber
\end{align}
If $\chi_\alpha\ge k>0$, then
\begin{equation}\label{inc:SGWR:v1:amperror}
 |\varepsilon_\alpha|\le 2\norm{\mathcal B_\alpha}/k.
\end{equation}
In particular, input/output alignment after initialization is not required.
\end{proposition}
\begin{proof}
The two Cartesian equations give
\[
 z_\alpha'=\frac\lambda2(1-c)z_\alpha
 +\frac12\big[(\mathcal B_\alpha^\top W_\alpha)_2
                    +(\mathcal B_\alpha U_\alpha)_2\big].
\]
Divide by positive $z_\alpha$ and differentiate $\log h_\alpha=2\log z_\alpha$.
The entrywise expression follows by grouping the two diagonal terms.
Each of the two defect contractions has magnitude at most
$\norm{\mathcal B_\alpha}\sqrt{m_\alpha}$, whereas
$z_\alpha\ge k\sqrt{m_\alpha}$. This proves the bound.
\end{proof}

\begin{corollary}[Growth-preserving cohort bound]\label{inc:SGWR:v1:ampcohort}
Suppose the same fixed set $C$ satisfies $\chi_\alpha\ge k>0$ and
$\norm{\mathcal B_\alpha(t)}\le d(t)$ for all its labels on
$[\tau_a,\tau_b]$, where $d$ is integrable. Let
$H_C(\tau)=\int_Ch_\alpha(\tau)\dd\lambda_0$.
Then
\begin{equation}\label{inc:SGWR:v1:ampbound}
 \boxed{N_C(\tau_b)\ge H_C(\tau_a)
 \exp\!\left[\lambda(\tau_b-\tau_a)
       -\lambda\int_{\tau_a}^{\tau_b}c(t)\dd t
       -\frac2k\int_{\tau_a}^{\tau_b}d(t)\dd t\right].}
\end{equation}
Here $H_C(\tau_a)\ge k^2N_C(\tau_a)$. At aligned isotropic initialization,
\begin{equation}\label{inc:SGWR:v1:ancestry}
 H_C(0)=S_0\int_C\sin^2\alpha\dd\lambda_0(\alpha).
\end{equation}
A sharper signed version replaces $-(2/k)\int d$ by the
$h_\alpha(\tau_a)$-weighted average of $\int\varepsilon_\alpha$.
\end{corollary}
\begin{proof}
Integrate \eqref{inc:SGWR:v1:amprate}, use the uniform lower bound on its
defect, and integrate over $C$. Since $h_\alpha\le m_\alpha$, this is a
lower bound for $N_C(\tau_b)$. The initial identity follows from
$U_{\alpha,2}(0)=W_{\alpha,2}(0)=\sqrt{S_0}\sin\alpha$.
The signed version follows by Jensen with the initial $h$ weights.
\end{proof}

\paragraph{The geometric deficit in the exact reference channel.}
For illustration, suppose \emph{exactly} that
$R=(\lambda/2)(1-c)e_2e_2^\top$ and that a label starts aligned at angle
$\alpha\in(\pi/2,\pi)$. This is an explicitly defined comparison channel,
not the assertion that the entire SIM residual has this form.
Let $\Lambda(t)=\lambda\int_{\tau_a}^{t}(1-c)\dd r$. Then
\[
 U_1=W_1=\sqrt{m_a}\cos\alpha,\qquad
 U_2=W_2=\sqrt{m_a}\sin\alpha\,e^{\Lambda/2},
\]
so
\begin{equation}\label{inc:SGWR:v1:channel}
 \log\frac{m(t)}{m_a}
 =\Lambda(t)+\log\big[\sin^2\alpha+e^{-\Lambda(t)}\cos^2\alpha\big]
 \ge\Lambda(t)+\log\sin^2\alpha.
\end{equation}
The initial angular misalignment costs only an endpoint constant for a
compact subarc, not a fixed fraction of the entire growth exponent.
Equivalently, $m\sin^2\theta=m_a\sin^2\alpha e^{\Lambda}$.
Equation \eqref{inc:SGWR:v1:amprate} is an exact perturbed two-angle version
of this observation. It does not prove its angular guards for the real flow.

\section{An exact Gaussian sufficient bound for the matrix defect}
The following estimate uses the full population output; it does not discard
background or replace a training cluster by its centroid.
Put $S=\int m_\alpha\dd\lambda_0$ and
\[
 \epsilon_2(c)=\norm{f(X)-cX_2e_2}_{L^2(P_2)}.
\]

\begin{proposition}[Interior-Q2 defect bound]\label{inc:SGWR:v1:gaussiandefect}
Suppose a receiving input direction satisfies $a_1\le-\eta<0$ and
$a_2\ge s_*>0$, and let $0\le c\le1$. For the defect in
\eqref{inc:SGWR:v1:defect},
\begin{align}
 \norm{\mathcal B_\alpha}\le \frac12\Big[&
 (1+S)(1+2q^2)\Phi(-\eta/q)+q^2\nonumber\\
 &+(\lambda+2q^2)\Phi(-\rho s_*/q)
 +\sqrt{\lambda+2q^2}\,\epsilon_2(c)\Big].
 \label{inc:SGWR:v1:gaussbound}
\end{align}
\end{proposition}
\begin{proof}
The balanced feature representation gives $\norm{f(X)}\le S|X|$.
For $X\sim\mathcal N(c_p,q^2I_2)$ and a unit direction $a$, direct
one-dimensional truncated-Gaussian integration yields, with $\ell=a^\top c_p$,
\[
 \E[|X|^2\1_{\{a^\top X>0\}}]
 =(|c_p|^2+2q^2)\Phi(\ell/q)+\ell q\phi(\ell/q).
\]
This follows by resolving $X$ into its projection along $a$ and an
independent orthogonal Gaussian, and integrating the first two moments
of a standard normal above $-\ell/q$. If $\ell<0$, its last term is
nonpositive. Consequently
\[
 \norm{R_1(a)}\le(1+S)(1+2q^2)\Phi(-\eta/q).
\]
For cluster 2 set $D_c=\operatorname{diag}(1,1-c)$ and
$e_c(X)=f(X)-cX_2e_2$. Then
\begin{align*}
 R_2(a)
 &=D_c\E_{P_2}[XX^\top]
   -D_c\E_{P_2}[XX^\top\1_{\{a^\top X\le0\}}]
   -\E_{P_2}[e_c(X)X^\top\1_{\{a^\top X>0\}}]\\
 &=\lambda(1-c)e_2e_2^\top+q^2D_c
   -D_c\E_{P_2}[XX^\top\1_{\{a^\top X\le0\}}]
   -\E_{P_2}[e_c(X)X^\top\1_{\{a^\top X>0\}}].
\end{align*}
Here $\norm{D_c}\le1$. Apply the same trace bound with receiver $-a$ to
the inactive moment. Cauchy--Schwarz bounds the last term by
$\sqrt{\lambda+2q^2}\epsilon_2(c)$. Average the two raw cluster fields,
retaining their mixture weight $1/2$, to obtain
\eqref{inc:SGWR:v1:gaussbound}.
\end{proof}

\begin{remark}[What must actually be controlled]
The norm estimate is only a sufficient budget; the signed contractions in
\eqref{inc:SGWR:v1:amprate} can be substantially better. For example,
negative $U_1,W_1$ paired with negative cross entries of $\mathcal B$ give
\emph{positive} growth corrections. Neither sign is assumed by the norm bound.
The scalar $\epsilon_2(c)$ includes the strong output, weak non-diagonal
output, and every uncharted label. It is not automatically $O(q)$ or
integrable merely because $N$ is small. Even a proved $O(q)$ bound over a
time of order $\log(1/S_0)$ initially gives an integrated cost of order
$q\log(1/S_0)$. An improved signed cancellation or a stronger integrated
bound is needed when that cost matters. Finally, the first tail bound ceases
to be exponentially small once $-a_1=O(q)$. The weak-chart transition needs
its own argument, not an extension of the interior-Q2 estimate by assertion.
\end{remark}

\section{Quantitative attraction in the frozen resident weak system}
This is a property of the defined scaled system with $M=1,N=0$ and the
strong atom fixed at $(a,a)$, where $a=\rho K(\rho a)>0$. It is not yet
capture under a moving population residual. Its equations are
\begin{equation}\label{inc:SGWR:v1:frozen}
 u'=F_\rho(u):=\tfrac12[\phi(u)-a\Phi(u)-\lambda u],
 \qquad v'=\tfrac12(a-\lambda v).
\end{equation}

\begin{proposition}[Frozen global attraction and a uniform contraction]
\label{inc:SGWR:v1:contraction}
For every $0<\rho<1$, the frozen system has a unique globally attracting
finite equilibrium $(u_*,a/\lambda)$. On the AN06 interval
$\rho\in[13/20,7/10]$ one has, for all real $u$,
\begin{equation}\label{inc:SGWR:v1:derivative}
 F_\rho'(u)\le-\frac{\lambda-1/4}{2}
 \le-\frac{69}{800}.
\end{equation}
Consequently
\[
 |u(t)-u_*|\le e^{-69(t-t_0)/800}|u(t_0)-u_*|,
 \quad
 |v(t)-a/\lambda|=e^{-\lambda(t-t_0)/2}|v(t_0)-a/\lambda|.
\]
For forced versions of \eqref{inc:SGWR:v1:frozen} with additive
$\epsilon_u,\epsilon_v$, the corresponding bounds have the additional
variation-of-constants terms
\[
 \int_{t_0}^t e^{-69(t-r)/800}|\epsilon_u(r)|\dd r,
 \qquad
 \int_{t_0}^t e^{-\lambda(t-r)/2}|\epsilon_v(r)|\dd r.
\]
\end{proposition}
\begin{proof}
The root lemma in the checkpoint shows that
$[\phi(u)-a\Phi(u)-\lambda u]/\Phi(u)$ is strictly decreasing from
positive infinity to negative infinity. Thus the scalar field is positive
below $u_*$ and negative above. Its solutions approach the unique root
monotonically from either side. The field has at most linear growth,
so no finite-time escape occurs. The $v$ equation is explicitly solvable.
For the uniform estimate,
\[
 F_\rho'(u)=\tfrac12[-(u+a)\phi(u)-\lambda].
\]
If $u\ge-a$, this is at most $-\lambda/2$. Otherwise put $t=-u>a$;
then $(t-a)\phi(t)\le t\phi(t)\le\phi(1)<1/4$.
The maximum follows by differentiating $t\phi(t)$ on $t\ge0$;
$\phi(1)<1/4$ follows from $\pi>3$ and $e>8/3$.
Since $\lambda\ge169/400$, \eqref{inc:SGWR:v1:derivative} follows.
The mean-value formula gives the Dini inequality for the difference from
the equilibrium. Integrating it, with or without the displayed forcing,
proves the remaining statements.
\end{proof}

\begin{remark}
The global assertion is global in the real coordinate of the \emph{frozen
scaled ODE}. Applying it to an original angle a fixed distance from $\pi/2$
would start at $u=O(1/q)$; validity of the scaled field there is not supplied
by a compact-chart expansion. Similarly, differences of $M$ from one,
nonzero $N$, strong mean/shape motion, and finite noise all enter the
forcing in an actual capture argument. They must be bounded, not assumed
absent because the frozen root is attracting.
\end{remark}

\section{Conditional use of the quoted global profile multiplier}
This section verifies algebra from the user's displayed formula only.
The cited file, its Theorem 5.1, and the definition of $b_c$ must be read
before marking its application established. In particular a one-sided
profile bound is not a two-sided bound on the full population distance.
Let $P_\rho=\Phi(\rho a_\rho)$ and $\alpha=\lambda P_\rho/2$.

\begin{proposition}[Global-multiplier exponent budget]
\label{inc:SGWR:v1:globalexponent}
Suppose at one common entry time
\[
 N_a\ge c_0S_0^{1-\lambda+b},\qquad
 \delta_a\le C_sS_0^{1/2-\alpha-b_s},\qquad 0<S_0<1,
\]
and suppose the profile quantity under consideration satisfies
\[
 \delta_b\le C_p
 \sqrt{\frac{M_a}{M_b}}\sqrt{\frac{N_b}{N_a}}
 \exp\!\left(\int_{\tau_a}^{\tau_b}b_c(t)\dd t\right)\delta_a
 +\epsilon_{\rm pass}.
\]
Assume $M_a/M_b\le C_M$, $N_b\le C_N$, and
\[
 \int_{\tau_a}^{\tau_b}b_c(t)\dd t\le B_p+b_p\log(1/S_0).
\]
Then
\begin{equation}\label{inc:SGWR:v1:globalpower}
 \delta_b\le C_pC_s\sqrt{C_MC_N/c_0}\,e^{B_p}
 S_0^{\eta_{\rm glob}-b/2-b_s-b_p}+\epsilon_{\rm pass},
 \qquad
 \eta_{\rm glob}=\frac\lambda2(1-P_\rho)>0.
\end{equation}
On $\rho\in[13/20,7/10]$, the AN06 bound $a_\rho<2/5$ gives
\begin{equation}\label{inc:SGWR:v1:rationalmargin}
 \eta_{\rm glob}>\frac{16393}{200000}.
\end{equation}
With no other exponent costs, the tolerable extra seed power is
$b<\lambda(1-P_\rho)$. The corresponding frozen-rate budget is
$b<(1-P_\rho)/P_\rho$.
\end{proposition}
\begin{proof}
Substitute the three bounds into the assumed profile estimate. The power
before the extra losses is
\[
 \frac12-\alpha-\frac{1-\lambda}{2}
 =\frac\lambda2-\frac{\lambda P_\rho}{2}.
\]
The constants retain their stated dependence on $q,\rho$.
For the rational estimate, $\phi(0)<2/5$ implies
\[
 P_\rho<\tfrac12+\tfrac25\tfrac7{10}\tfrac25=\frac{153}{250},
 \qquad
 \frac\lambda2(1-P_\rho)>
 \frac{169}{800}\frac{97}{250}=\frac{16393}{200000}.
\]
The extra seed power costs $b/2$ here. Under amplification
$N_a^{-\alpha/\lambda}$ it instead costs $bP_\rho/2$, against baseline
$(1-P_\rho)/2$, giving the other budget.
\end{proof}

\paragraph{Why an integrated center bound is still a premise.}
For the exact leading logistic masses on an interval ending at $N_b<1$,
\begin{equation}\label{inc:SGWR:v1:logisticintegrals}
 \int_{\tau_a}^{\tau_b}N\dd\tau
 =\frac1\lambda\log\frac{1-N_a}{1-N_b},\qquad
 \int_{\tau_a}^{\tau_b}(1-M)\dd\tau=\log\frac{M_b}{M_a}.
\end{equation}
Thus $\int N\le(\log2)/\lambda$ up to half-mass. If, for example, a
proved estimate gave
$(b_c)_+\le C[N+(1-M)]+e_c$ with $\int|e_c|\le B$ and $M_a$ bounded
away from zero, the exponential multiplier would have an $O(1)$ exponent.
A similar statement holds for a quantity controlling $b_c$.
Vanishing at the resident point, or pointwise smallness, alone does not
supply this integral estimate uniformly as $N_a\to0$. The full $b_c$ or
$U_c$ definition and its dynamical bound must be supplied by the missing
source and subsequent analysis.

\section{Precise next initialized target}
The signed-growth change alters the main retention interface, but it does
not supply the final seed. A useful next lemma must specify a fixed set of
initial labels, prove its angular guards through the bulk stage, quantify
its ancestry amplitude $H_C$, and bound the \emph{signed} growth defect.
It must then transfer a quantified mass fraction into a bounded weak chart
while controlling the actual moving strong residual and all background
labels. The frozen contraction proposition supplies a terminal comparison
only after its forcing and domain assumptions have been justified.

If $N_C(\tau_a)\ge cS_0^\gamma$ and
$\tau_b-\tau_a=B_t\log(1/S_0)+O(1)$, a signed lower growth estimate
with logarithmic loss $b\log(1/S_0)$ gives a seed lower bound with power
$\gamma-\lambda B_t+b$. One must not count the same amplification both in
$N_C(\tau_a)$ and in the elapsed time. At finite noise the clock coefficient
$1/(1+2q^2)$, when proved for the relevant strong clock, changes this power;
its effect over long delays belongs to the finite-noise budget.

The interior-Q2 bulk route and the earlier backward-selected reservoir
route use different label sets and parameter regimes. A conservative
$S_0^{17/16}$ reservoir bound does not become a captured
$S_0^{1-\lambda}$ bound without an amplification-and-capture argument.
Nor does the user's direct masked seed measurement prove that the same
labels are retained.

A total signed deficit $O(1)$ loses only a prefactor. A deficit
$o(\log(1/S_0))$ permits arbitrarily small power losses but can still have
an unbounded prefactor. Either can be sufficient if the profile exponent
margin and all $q$-dependent factors remain controlled. The fitted
retention cost near $0.006$ is motivation, not a premise of these results.

\paragraph{What is not complete.}
The initialized angular guards; the original-flow weak seed at the desired
scale; nonlinear center and two-sided shape control; moving chart capture;
background and exact probe-rate estimates; canonical finite-parameter
coverage; and initialized Theorems A and B remain open. The missing
ancestry/tail theorem must be re-uploaded before its precise hypotheses can
be incorporated. No new numerical experiments were needed for any displayed
proof in this increment.
\end{document}
